\documentclass{beamer}
\usepackage{beamerthemesplit}
\setbeamertemplate{items}[ball]
\usetheme[height=7mm]{Rochester}
\hypersetup{pdfpagemode=FullScreen}
\usepackage{geometry}
\usepackage{graphicx}
\usepackage{amssymb}
\usepackage{amsthm}
\usepackage{color}
\usepackage{mathrsfs}
\usepackage[all]{xy}
\usepackage{listings}
\lstset{frame=single,
        language=SAS,
        basicstyle=\tiny\tt,
        keywordstyle=\color{black},
}
\usepackage[latin1]{inputenc}
\usepackage{pgf,pgfarrows,pgfnodes}
\beamertemplateballitem
\usepackage{hyperref}
\setbeamertemplate{footline}[frame number]


\hypersetup{
  colorlinks = true,
  allcolors= blue,
  pdfauthor = {KB Christensen},
  pdfkeywords = {Scale validation},
  pdftitle = {Scale validation},
  pdfsubject = {Scale validation},
  pdfpagemode = UseNone
}

\title{(iv)  Local independence:\ $X_i\perp X_j|\Theta$}
\author{Karl Bang Christensen \\ \href{mailto:kach@sund.ku.dk}{mailto:kach@sund.ku.dk}}
\date{Sect. of Biostatistics\\Dept. of Public Health\\Univ. of Copenhagen\\ \url{http://publicifsv.sund.ku.dk/~kach/scaleval/}}

\begin{document}

\maketitle

\begin{frame}\frametitle{Model framework, ctd.}
Latent variable $\Theta$, items $(X_i)_{i\in I}$, exogenous variables $Y$
$$
\xymatrix{
            &                    &X_1\\
Y\ar@{<->}[rd]&                    &   &X_2\\
            &\Theta\ar@{->}[uur]
                  \ar@{->}[urr]
                  \ar@{->}[rr]
                  \ar@{->}[rd]  &   &X_3\\
            &                   &X_4\\
}
$$
\begin{itemize}
\item[]
\end{itemize}
\end{frame}


\begin{frame}\frametitle{Model framework, Assumptions}
Latent variable $\Theta$, items $(X_i)_{i\in I}$, exogenous variables $Y$
$$
\xymatrix{
            &                    &X_1\\
Y\ar@{<->}[rd]&                    &   &X_2\\
            &\Theta\ar@{->}[uur]
                  \ar@{->}[urr]
                  \ar@{->}[rr]
                  \ar@{->}[rd]  &   &X_3\\
            &                   &X_4\\
}
$$
\begin{itemize}
\item[(iv)] Local independence:\ $X_i\perp X_j|\Theta$
\end{itemize}
\end{frame}


\begin{frame}\frametitle{Model framework, Assumptions}
Latent variable $\Theta$, items $(X_i)_{i\in I}$, exogenous variables $Y$
$$
\xymatrix{
            &                    &X_1\ar@{-}[dr]^{?}\\
Y\ar@{<->}[rd]&                    &   &X_2\\
            &\Theta\ar@{->}[uur]
                  \ar@{->}[urr]
                  \ar@{->}[rr]
                  \ar@{->}[rd]  &   &X_3\\
            &                   &X_4\\
}
$$
\begin{itemize}
\item[(iv)] Local independence:\ $X_i\perp X_j|\Theta$
\end{itemize}
\end{frame}


\begin{frame}\frametitle{Assumptions are testable }
Latent variable $\Theta$, items $(X_i)_{i\in I}$, exogenous variables $Y$
\begin{itemize}
\item Scale is valid $\Rightarrow$ Assumptions (i)-(iv) are met
\item Assumptions (i)-(iv) are \textit{not} met $\Rightarrow$ Scale is \textit{not} valid
\end{itemize}
\end{frame}

\begin{frame}\frametitle{Partial correlation}
measures the degree of association between two random variables, with the effect of a set of controlling random variables removed
$$
\hat{\rho}_{XY\cdot\mathbf{Z}}=\frac{N\sum_{i=1}^N r_{X,i}r_{Y,i}-\sum_{i=1}^N r_{X,i}\sum_{i=1}^N r_{Y,i}}
{\sqrt{N\sum_{i=1}^N r_{X,i}^2-\left(\sum_{i=1}^N r_{X,i}\right)^2}~\sqrt{N\sum_{i=1}^N r_{Y,i}^2-\left(\sum_{i=1}^N r_{Y,i}\right)^2}}
$$
\end{frame}


\begin{frame}\frametitle{Partial correlation:\ example}
{\small
\begin{itemize}
\item Wechsler Adult Intelligence Scale (WAIS). Correlation between sub-scales \underline{C}omprehension, \underline{A}rithmetic ability, \underline{V}ocabulary
\item[]\begin{tabular}{cc}
                                            &\begin{tabular}{ccc}C&A&V\\\end{tabular}\\
    \begin{tabular}{c}C\\A\\V\\\end{tabular}&\begin{tabular}{|ccc|}\hline
                                             1&0.49&0.73\\
                                              &1   &0.59\\
                                              &    &1\\\hline
                                             \end{tabular}                           \\
    \end{tabular}
\item Correlation $\rho_{CA}=0.49$. Partial correlation $\rho_{CA|V}=0.11$
\item With the effect of vocabulary removed, the correlation between comprehension and arithmetic ability disappears
\item WAIS in a sample of subjects who are homogeneous wrt. vocabulary: correlation between comprehension scores arithmetic scores small.
\item[]
\item .. 'explains the association'.
\end{itemize}
}
\end{frame}


\begin{frame}\frametitle{Construct validity}
latent variable is the only variable that influences item responses
observed item responses can provide measurement of the value of the latent variable.

\xymatrix@=1.1em{
     &X_1\;&                                                                     &EXO_1\ar[ddl]        \\
X_2\;&     &                                                                     &                     \\
           &     &\ar@/_/[uul]\ar@/_/[ull]\ar@/^/[dll]\ar@/^/[ddl] \;\;\theta\;\;&                    &\\
X_3\;&     &                                                                     &EXO_3\ar[uuu]\ar[ul] \\
     &X_4\;&                                                                     &
}
\end{frame}


\begin{frame}{Differential item functioning (DIF)}
Item responses influenced by other variables - measurement will be confounded

\xymatrix@=1.1em{
&     &X_1\;            &                                                               &EXO_1\ar[ddl]        \\
&X_2\;&                 &                                                               &                     \\
&     &                 &\ar@/_/[uul]\ar@/_/[ull]\ar@/^/[dll]\ar@/^/[ddl] \;\;\theta\;\;&                    &\\
&X_3\;&                 &                                                               &EXO_3\ar[uuu]\ar[ul] \\
&     &X_4\;\ar@{-}[rru]&                                                               &
}

Test:\ conditional independence
$$
X_4\perp EXO_3|R
$$
$R=X_1+X_2+X_3+X_4$.
\end{frame}


\begin{frame}{Local dependence}
\xymatrix@=1.1em{
&                   &X_1\;&                                                               &EXO_1\ar[ddl]       \\
&X_2\;              &     &                                                               &                    \\
&                   &     &\ar@/_/[uul]\ar@/_/[ull]\ar@/^/[dll]\ar@/^/[ddl] \;\;\theta\;\;&                    \\
&X_3\ar@{-}[rd]\;&        &                                                               &EXO_3\ar[uuu]\ar[ul]\\
&                   &X_4\;&                                                               &
}
Test:\ ?
\end{frame}


\begin{frame}{Local dependence as DIF}
treat $X_4$ as exogenous variable. Test for DIF

\xymatrix{
&        &X_1\;\ar@{-}[rr]&                                           &EXO_1\ar[ddl]                  \\
&X_2\;&        &                                                      &                             \\
&        &        &\ar@/_/[uul]\ar@/_/[ull]\ar@/^/[dll] \;\;\theta\;\;&                   &\\
&X_3\;&        &                                                      &EXO_3\ar[uuu]\ar[ul]           \\
&        &     &                                                      &X_4\;\ar@{-}[lllu]\ar[uul]
}
Test:\ conditional independence
$$
X_3\perp X_4|R_4
$$
$R_4=X_1+X_2+X_3$.
\end{frame}


\begin{frame}{Check local independence}
We can test local independence using
$$
X_i\perp X_k|R_i
$$
and
$$
X_i\perp X_k|R_k
$$
based on test proposed by Tjur\footnote{Tjur. Scandinavian Journal of Statistics 9, 1982, 23-30.}. Use
\begin{itemize}
\item Partial correlation coefficents
\item Mantel-H{\"a}entzel
\item Logistic regression
\end{itemize}
\end{frame}


\begin{frame}{Check local independence in the SCQOL data}
\lstinputlisting[linerange=3-15]{(iv)_Local_independence.sas}
\end{frame}


\begin{frame}{Check local independence in the SCQOL data}
\begin{center}
\begin{tabular}{|rrrrrrrrr|}
\hline
 1.00 & 0.01 &-0.27 &-0.08 & 0.24 & 0.04 & 0.05 &-0.16 & 0.04 \\
 0.03 & 1.00 &-0.04 & 0.16 &-0.08 &-0.05 &-0.03 &-0.02 & 0.06 \\
-0.19 & 0.05 & 1.00 & 0.11 &-0.05 & 0.06 & 0.04 & 0.62 &-0.05 \\
-0.03 & 0.20 & 0.06 & 1.00 &-0.13 &-0.04 & 0.16 & 0.05 & 0.05 \\
 0.22 &-0.13 &-0.15 &-0.21 & 1.00 & 0.22 &-0.14 &-0.10 &-0.02 \\
 0.07 &-0.03 & 0.01 &-0.05 & 0.27 & 1.00 & 0.10 &-0.15 & 0.00 \\
 0.05 &-0.06 &-0.06 & 0.10 &-0.12 & 0.07 & 1.00 &-0.13 & 0.00 \\
-0.18 &-0.06 & 0.56 &-0.03 &-0.08 &-0.20 &-0.15 & 1.00 &-0.10 \\
 0.03 & 0.02 &-0.17 &-0.03 & 0.01 &-0.05 &-0.01 &-0.09 & 1.00 \\
\hline
\end{tabular}
\end{center}
\end{frame}


\begin{frame}{Check local independence in the SCQOL data}
\begin{center}
\begin{tabular}{|rrrrrrrrr|}
\hline
 1.00 & 0.01 &-0.27 &-0.08 &{\bf 0.24} & 0.04 & 0.05 &-0.16 & 0.04 \\
 0.03 & 1.00 &-0.04 & 0.16 &-0.08 &-0.05 &-0.03 &-0.02 & 0.06 \\
-0.19 & 0.05 & 1.00 & 0.11 &-0.05 & 0.06 & 0.04 & 0.62 &-0.05 \\
-0.03 & 0.20 & 0.06 & 1.00 &-0.13 &-0.04 & {\bf 0.16} & 0.05 & 0.05 \\
{\bf  0.22} &-0.13 &-0.15 &-0.21 & 1.00 & {\bf 0.22} &-0.14 &-0.10 &-0.02 \\
 0.07 &-0.03 & 0.01 &-0.05 & {\bf 0.27} & 1.00 & 0.10 &-0.15 & 0.00 \\
 0.05 &-0.06 &-0.06 & 0.10 &-0.12 & 0.07 & 1.00 &-0.13 & 0.00 \\
-0.18 &-0.06 & {\bf 0.56} &-0.03 &-0.08 &-0.20 &-0.15 & 1.00 &-0.10 \\
 0.03 & 0.02 &-0.17 &-0.03 & 0.01 &-0.05 &-0.01 &-0.09 & 1.00 \\
\hline
\end{tabular}
\end{center}
\end{frame}


\begin{frame}{Check local independence in the COPD data}
\begin{enumerate}
    \item From the course homepage download the sample SAS code \url{http://biostat.ku.dk/~kach/scaleval/exerc7.sas}
    \item Use the {\%}tjur macro to compute partial correlations
    \item Test local independence for item pairs with high partial correlations
\end{enumerate}
\end{frame}


\end{document}
